\documentclass[10pt,a4paper]{amsart}
\usepackage[margin=19mm]{geometry}
\usepackage{amsmath,amssymb,mathtools}
\usepackage[scr=rsfs,cal=euler]{mathalfa}
\usepackage{xfrac}
\usepackage[unicode,hidelinks,bookmarks=false]{hyperref}
\newcommand{\ie}{i.e.\ }
\newcommand{\cf}{cf.\ }
\newcommand{\resp}{respectively\ }
\newcommand{\Subsection}{\subsection}
\providecommand{\nobreakdash}{\nobreak\hbox{-}}
\setlength{\emergencystretch}{2em}
\multlinegap0pt
\usepackage{amssymb}
\newtheorem{theorem}{Theorem}
\newcommand{\C}{{\mathbb{C}}}
\renewcommand{\P}{{\mathbb{P}}}
\title{Meromorphic functions bi-weighted weakly sharing pairs of small functions}
\author{Si Duc Quang, Phung Nguyen Ngoc Anh}
\date{Source: arXiv:2605.26987v1 (CC BY 4.0)}
\begin{document}
\maketitle

\noindent\emph{Source and licence.} Meromorphic functions bi-weighted weakly sharing pairs of small functions. Version arXiv:2605.26987v1 (CC BY 4.0), distributed by arXiv under the Creative Commons Attribution 4.0 International licence, http://creativecommons.org/licenses/by/4.0/, as recorded in the arXiv licence metadata for this exact version and shown on the abstract page https://arxiv.org/abs/2605.26987v1. Full licence notice: \texttt{licenses/arxiv-2605.26987-CC-BY-4.0.txt}.\par\smallskip

\noindent\textit{Editorial scope.} Counted unit: the article's theorem 1.2 (source0.tex lines 189--196). The article's complete original proof of it is delivered with it (source0.tex lines 496--592, one \texttt{proof} environment of 97 lines, which the article places away from the statement). No mathematics is written, repaired or re-derived: the statement and the proof are the article's own lines, in the article's own notation, and no retained line is substituted or re-flowed.  Retained uncounted as the source-local context the proof needs: lines 303--305.  Nothing was omitted from any retained span, so the package carries no omission record.  No retained line contains a citation, so the wrapper carries no bibliography and no bibliography entry.  No retained line contains a figure environment, an image-inclusion command or drawing code, so no image is delivered and the package declares no asset dependency.  Editorial formatting only, in addition: an AMS \texttt{amsart} preamble that restates the article's own declarations and macros used by the retained lines, namely the environments \texttt{theorem} on separate counters, together with the standard packages those lines need.  The retained environments print in this document as Theorem 1 in order of appearance, whereas the article prints the same items with the numbering of its own full document.  The complete native source of the article as distributed in the arXiv e-print for this exact version is bundled unchanged as \texttt{sources/201459/source0.tex}.
\begin{align}\label{3.5}
\|\ \Delta\overline N_{(n_1+1,k)}(r,\nu^{a_1}_f)\le \frac{2(n_2+1)(n_3+1)}{k+1}T(r)+S(r),
\end{align}

\begin{theorem}\label{1.2}
Let $f$ and $g$ be non-constant meromorphic functions and $a_i,b_i\ (i=1,2,3,4)$ $ (a_i\ne a_j, b_i\ne b_j, i\ne j)$ be small functions (with respect to $f$ and $g$).  Let $n_1,n_2,n_3$ and $k$ be positive integers such that $\Delta:=n_1n_2n_3-n_1-n_2-n_3-2>0$ and there exists a positive integer $n$ satisfying
$$\frac{48\sum_{1\le i<j\le 3}(n_i+1)(n_j+1)+72\Delta}{\Delta(k+1)}+\frac{30\lambda_n}{k+1}+\frac{88}{n+2}<1,$$
where $\lambda_n=2^{n^2+2n+3}\left(\frac{2\sum_{1\le i<j\le 3}(n_i+1)(n_j+1)}{\Delta(k+1)}+\frac{3}{k+1}\right)$.
If $f$ and $g$ weakly share the pair $(a_i,b_i)$ with bi-weight $(n_i,k)$ for every $i=1,2,3$ and satisfy a condition that
$$E^S_{1)}(a_4,f)=E^S_{1)}(b_4,g),E^S_{2)}(a_4,f)\subset E^S_{2}(b_4,g),E^S_{2)}(b_4,g)\subset E^S_{2}(a_4,f)$$
for a discrete subset $S$ of with $\| N(r,\nu_S)=o(T(r,f)+T(r,g))$, then $f$ and $g$ are quasi-M\"{o}bius transformation of each other.
\end{theorem}

\begin{proof}[Proof of Theorem \ref{1.2}]
By using a quasi-M\"{o}bius transformation if necessary, we may suppose that $a_1=b_1=0,a_2=b_2=\infty,a_3=b_3=1,a_4=a,b_4=b$ and $a,b\not\in\{0,1,\infty\}$. Suppose by contrary that $f$ is not a quasi-M\"{o}bius transformation of $g$. We set $T(r)=T(r,f)+T(r,g)$ and $S(r)=S(r,f)+S(r,g).$ 
We define divisors 
$$\mu_i(z)=\min\{1,|\nu^0_{f-a_i}(z)-\nu^0_{g-b_i}(z)|\}\text{ for }i=1,2,3\text{ and }\mu_4(z)=|\nu^0_{f-a_4}(z)-\nu^0_{g-b_4}(z)|.$$
By the supposition and by (\ref{3.5}), we have
\begin{align*}
\|\ \Delta\sum_{i=1}^3N(r,\mu_i)&\le\Delta\sum_{i=1}^3\left(\overline N_{(n_i+1,k)}(r,\nu^{a_i}_f)+\overline N_{(k+1}(r,\nu^{a_i}_f)+\overline N_{(k+1}(r,\nu^{b_i}_g)\right)+S(r)\\
&\le \frac{2\sum_{1\le i<j\le 3}(n_i+1)(n_j+1)+3\Delta}{k+1}T(r)+S(r).
\end{align*}
Also, by the main theorem, we have
$$\|\ N(r,\mu_4)\le N_{(3}(r,\nu^{a_4}_f)+N_{(3}(r,\nu^{b_4}_g)+S(r)\le 2\left(\frac{2\lambda_n}{k+1}+\frac{6}{n+2}\right)T(r)+S(r),$$
where
$$\lambda_n=2^{n^2+2n+3}\left(\frac{2\sum_{1\le i<j\le 3}(n_i+1)(n_j+1)}{\Delta(k+1)}+\frac{3}{k+1}\right)$$
for a positive integer $n$ (chosen later). Combining the above two inequalities, we obtain
\begin{align}\label{4.2}
\begin{split}
4&\sum_{i=1}^3N(r,\mu_i)+N(r,\mu_4)\\
&\le \left(\frac{8\sum_{1\le i<j\le 3}(n_i+1)(n_j+1)+12\Delta}{\Delta(k+1)}+\frac{4\lambda_n}{k+1}+\frac{12}{n+2}\right)T(r)+S(r).
\end{split}
\end{align}
We write $f=\frac{f_0}{f_1}$ (resp. $g=\frac{g_0}{g_1}$), where $f_0,f_1$ (resp. $g_0,g_1$) are holomorphic function without common zeros. Take meromorphic functions $h_1,h_2,h_3,h_4$ such that
$$f_0=h_1 g_0;f_1=h_2 g_1;f_0-f_1=h_3(g_0-g_1);f_0-af_1=h_4(g_0-bg_1).$$
It is easy to see that 
$$\nu^0_{h_i}+\nu^\infty_{h_i}=|\nu^{a_i}_f-\nu^{b_i}_g|\ (1\le i\le 4)$$
outside a discrete subset $S$ of $\C$ with $\|\ N(r,\nu_S)=S(r)$.
From the definition of functions $h_i\ (1\le i\le 4)$, we have 
\begin{align*}
\mathrm{det}\left (
\begin{array}{cccc}
1&0&-h_1&0\\ 
0&1&0&-h_2\\
1&-1&-h_3&h_3\\
1&-a&-h_4&bh_4
\end{array}\right )=0.
\end{align*}
Then
\begin{align*}
%\label{3.12}
(1-a)h_1h_2-h_1h_3+bh_1h_4+a h_2h_3-h_2h_4+(1-b) h_3h_4=0.
\end{align*}
We define the following functions:
\begin{align*}
h_{1,2}&=(1-a)h_1h_2,\ h_{1,3}=-h_1h_3,\ h_{1,4}=b h_1h_4,\\
h_{2,3}&=a h_2h_3,\ h_{2,4}=-h_2h_4,\ h_{3,4}=(1-b) h_3h_4.
\end{align*}
Then we have 
$$ \sum_{1\le i<j\le 4}h_{i,j}=0. $$
Let $d$ be a meromorphic function on $\C$ such that $dh_{i,j}\ (1\le i<j\le 4)$ are all holomorphic functions on $\C$ without common zero. Then it is easy to see that
\begin{align*}
 \sum_{1\le i<j\le 4}N^{[4]}(r,dh_{i,j})\le& 3\sum_{i=1}^4(N^{[4]}(r,\nu^0_{h_i})+N^{[4]}(r,\nu^\infty_{h_i}))+S(r)\\ 
\le& 12\sum_{i=1}^3N(r,\mu_i)+3N(r,\mu_4)+S(r).
\end{align*}
Let $\mathcal I=\{(i,j)|1\le i<j\le 4\}$. For $I=(i,j)\in\mathcal I$, we denote $h_{i,j}$ by $h_{I}$.
Take $I_0\in\mathcal I$ arbitrarily. Then 
$$dh_{I_0}=-\sum_{I\ne I_0}dh_I.$$
Denote by $t$ the minimum number satisfying the following: There exist $t$ elements $I_1$,..., $I_t \in\mathcal I$ and $t$ nonzero constants $b_v\in\C\ (1\le v\le t)$ such that $dh_{I_0}=\sum_{v=1}^tb_vdh_{I_v}.$

By the minimality of $t$, the family $\{dh_{I_1},...,dh_{I_t}\}$ is linearly independent over $\C$.

Suppose that $t\ge 2$. Consider the linearly non-degenerate holomorphic mapping $h:\C \to \P^{t-1}(\C)$ with the representation $h=(d h_{I_1}:...:d h_{I_t})$. By the second main theorem in Nevanlinna theory for hyperplanes, we have
\begin{align}\label{4.4}
\begin{split}
\|\ T_h(r)&\le \sum_{v=1}^t N^{[t-1]}_{dh_{I_v}}(r)+N_{dh_{I_0}}^{[t-1]}(r)+S(r)\\
&\le \sum_{v=1}^6 N^{[4]}_{dh_{I_v}}(r)+S(r)\ \  [\text{since $t\le 5$}]\\
&\le 12\sum_{i=1}^3N(r,\mu_i)+3N(r,\mu_4)+S(r).
\end{split}
\end{align}
On the other hand, there exists two element, say $I_i$ and $I_j$, such that $4\not\in (I_i\setminus I_j)\cup (I_j\setminus I_i)$. Therefore, by the main theorem we have
\begin{align}\label{4.5}
\begin{split}
\|\ T_h(r)&\ge T(r,\frac{h_{I_i}}{h_{I_j}})+S(r)=T(r,\frac{h_{I_i}}{h_{I_j}}-\frac{\tilde h_{I_i}}{\tilde h_{I_j}})+S(r)\\
&\ge N(r,\min\{\nu^0_{f-a},\nu^0_{g-b}\})+S(r)\\
&\ge \frac{1}{2}\left(N(r,\nu^a_{f})+N(r,\nu^a_{f})-N_{(3}(r,\nu^a_{f})-N_{(3}(r,\nu^a_{f})\right)+S(r)\\
&\ge \frac{1}{2}\left(1-\frac{6\lambda_n}{k+1}-\frac{16}{n+2}\right)T(r)+S(r),
\end{split}
\end{align}
where $\tilde h_{I_i},\tilde h_{I_j}$ are obtained from $h_{I_i},h_{I_j}$ by substituting $a,b$ into $f,g$ respectively.
From (\ref{4.2}), (\ref{4.4}) and (\ref{4.5}), we have
\begin{align*}
&\biggl\|\  \frac{1}{2}\left(1-\frac{6\lambda_n}{k+1}-\frac{16}{n+2}\right)T(r)\\ 
& \le 3\left(\frac{8\sum_{1\le i<j\le 3}(n_i+1)(n_j+1)+12\Delta}{\Delta(k+1)}+\frac{4\lambda_n}{k+1}+\frac{12}{n+2}\right)T(r)+S(r).
\end{align*}
Letting $r\longrightarrow +\infty$, we get 
$$1\le \frac{48\sum_{1\le i<j\le 3}(n_i+1)(n_j+1)+72\Delta}{\Delta(k+1)}+\frac{30\lambda_n}{k+1}+\frac{88}{n+2}.$$
This is a contradiction.

Therefore, we must have $t=1$, i.e., there is $I_1\in \mathcal {I}\setminus \{I_0\}$ such that  $\dfrac {h_{I_0}}{h_{I_1}}\in\C\setminus\{0\}$. Hence, for every $I\in\mathcal J$, there exists $J\in\mathcal I\setminus\{I\}$ such that $\dfrac{h_I}{h_J}\in\C\setminus\{0\}$. We consider the following three cases:

\textbf{Case a:} There are $I=(s,t), J=(s,\ell)$ such that $\dfrac{h_I}{h_J}\in\C\setminus\{0\}$. Then $h_t=\alpha h_\ell$ with a small meromorphic function (w.r.t $f$ and $g$) $\alpha$. Therefore, $f$ is a quasi-M\"{o}bius  transformation of $g$. This contradicts the supposition.

\textbf{Case b:} There are $I=(t,s), J=(\ell,s)$ such that $\dfrac{h_I}{h_J}\in\C\setminus\{0\}$.  Similarly as the above case, $f$ is a quasi-M\"{o}bius  transformation of $g$ and we arrive at the contradiction. 

\textbf{Case c:} There are small functions $\alpha,\beta$ such that $h_{1,2}=\alpha h_{3,4}\text{ and } h_{1,3}=\beta h_{2,4}.$
Then $\left(\dfrac{h_1}{h_4}\right)^{2}$ is a small function, and so is $\dfrac{h_1}{h_4}$. Hence $f$ is a quasi-M\"{o}bius transformation of $g$. This is a contradiction.

From the above two cases, we always get a contradiction. Hence $f$ is a quasi-M\"{o}bius transformation of $g$. 
\end{proof}

\end{document}
